% questions.tex —— 1992 年·三南卷 · 题目内容片段（无答案）
% 由 tools/generate_projects.py 依来源总账自 1992/1992_数学真题_三南_试题.md 生成，勿手改（改源 MD 后重跑）


\gkpart{数学}

\begin{gktip}
\textbf{来源说明}：依据固定来源 PDF 的原生文字层逐题转录，并对照 400 dpi 页面渲染逐题核对。源为 \textbf{XeTeX 重排件（非扫描）}——全矢量、文字层可读、嵌入 LM／CM 系与 STSong／STHeiti 字体，含 CMMI10／CMSY10。单页 A3 横向三栏（页面 1190.55×841.89 pt），页脚仅印印刷页码 \textbf{44}；按原页三栏次序为 1—10、11—23、24—28；\textbf{源页无大题标题}，题号 1–28 连续；\textbf{无答案页，无题图}。

\textbf{标点与字样约定}：① 题号源内为半角句点，本库统一作全角「．」；② 源内句末为实心点「．」，按全库约定归一为「。」，其余中文标点归一全角；③ 选项标号源为半角「(A)」，排印统一作全角「（A）」；④ 分数按源一律竖排真分式，排印用 \(\dfrac{a}{b}\)；⑤ 根号、上下标照源；⑥ 虚数单位 \(\mathrm{i}\) 为正体，实数集 \(\mathbb R\)、整数集 \(\mathbb Z\)、自然数集 \(\mathbb N^*\) 按源。

\textbf{存疑与登记}（源内异常照录，不推测补正）：① \textbf{第 1 题}用 \(\arg z\) 而未声明主值区间（主值与全体辐角两种约定可给出不同选项），题面照录，作答约定见答案卷。② \textbf{第 11 题选项 B} 印作 \(2(2-\sqrt3\pi)\)——\(\pi\) 落在括号内，与其余选项量纲不一致（疑为 \(2(2-\sqrt3)\pi\) 的排印移位）；照录原样。③ \textbf{第 15 题}选项 A、D 均称「\(A<B\)」是「\(\tan A<\tan B\)」的充分条件，其中 D 项印作「充分条件，也不是必要条件」，漏「既不」二字。题面照录，判定见答案卷。④ \textbf{第 24 题}指数 \(2x-2\) 与 \(x-1\) 已经 400 dpi 复核；⑤ \textbf{第 27 题}截线长度 \(4\sqrt{10}\) 已经复核。

\textbf{图示}：源页无题图（第 5、23 题的立体图形仅以文字描述）。
\end{gktip}

\begin{question}
\gkti{1}设函数 $z=\mathrm{i}^2+\sqrt3\,\mathrm{i}$，那么 $\arg z$ 是\mbox{（　　）}。

（A）$\dfrac{5\pi}{6}$　　（B）$\dfrac{\pi}{3}$　　（C）$\dfrac{2\pi}{3}$　　（D）$-\dfrac{4\pi}{3}$

\end{question}
\begin{question}
\gkti{2}如果等边圆柱（即底面直径与母线相等的圆柱）的体积是 $16\pi\,\mathrm{cm}^3$，那么它的底面半径等于\mbox{（　　）}。

（A）$4\sqrt[3]{2}\,\mathrm{cm}$　　（B）$4\,\mathrm{cm}$　　（C）$2\sqrt[3]{2}\,\mathrm{cm}$　　（D）$2\,\mathrm{cm}$

\end{question}
\begin{question}
\gkti{3}$\dfrac{\arcsin\dfrac{\sqrt3}{2}-\arccos\left(-\dfrac12\right)}{\arctan(-\sqrt3)}$ 的值等于\mbox{（　　）}。

（A）$1$　　（B）$0$　　（C）$-\dfrac25$　　（D）$-\dfrac65$

\end{question}
\begin{question}
\gkti{4}函数 $y=\log_{1/2}(1-x)\ (x<1)$ 的反函数是\mbox{（　　）}。

（A）$y=1+2^{-x}\ (x\in\mathbb R)$　　（B）$y=1-2^{-x}\ (x\in\mathbb R)$

（C）$y=1+2^x\ (x\in\mathbb R)$　　（D）$y=1-2^x\ (x\in\mathbb R)$

\end{question}
\begin{question}
\gkti{5}在长方体 $ABCD-A_1B_1C_1D_1$ 中，如果 $AB=BC=a$，$AA_1=2a$，那么点 $A$ 到直线 $A_1C$ 的距离等于\mbox{（　　）}。

（A）$\dfrac{2\sqrt6}{3}a$　　（B）$\dfrac{3\sqrt6}{2}a$　　（C）$\dfrac{2\sqrt3}{3}a$　　（D）$\dfrac{\sqrt6}{3}a$

\end{question}
\begin{question}
\gkti{6}函数 $y=\sin x\cos x+\sqrt3\cos^2x-\dfrac{\sqrt3}{2}$ 的最小正周期等于\mbox{（　　）}。

（A）$\pi$　　（B）$2\pi$　　（C）$\dfrac\pi4$　　（D）$\dfrac\pi2$

\end{question}
\begin{question}
\gkti{7}有一个椭圆，它的极坐标方程是\mbox{（　　）}。

（A）$\rho=\dfrac5{\sqrt3-2\cos\theta}$　　（B）$\rho=\dfrac5{\sqrt3-\sqrt3\cos\theta}$

（C）$\rho=\dfrac{2-\sqrt3\cos\theta}{5}$　　（D）$\rho=\dfrac5{2-\sqrt3\cos\theta}$

\end{question}
\begin{question}
\gkti{8}不等式 $\left|\sqrt{x-2}-3\right|<1$ 的解集是\mbox{（　　）}。

（A）$\{x\mid 5<x<16\}$　　（B）$\{x\mid 6<x<18\}$　　（C）$\{x\mid 7<x<20\}$　　（D）$\{x\mid 8<x<22\}$

\end{question}
\begin{question}
\gkti{9}设等差数列 $\{a_n\}$ 的公差是 $d$，如果它的前 $n$ 项和 $S_n=-n^2$，那么\mbox{（　　）}。

（A）$a_n=2n-1,\ d=-2$　　（B）$a_n=2n-1,\ d=2$

（C）$a_n=-2n+1,\ d=-2$　　（D）$a_n=-2n+1,\ d=2$

\end{question}
\begin{question}
\gkti{10}方程 $\cos2x=3\cos x+1$ 的解集是\mbox{（　　）}。

（A）$\{x\mid x=2k\pi\pm\dfrac{2\pi}{3},\ k\in\mathbb Z\}$　　（B）$\{x\mid x=k\pi\pm\dfrac\pi3,\ k\in\mathbb Z\}$

（C）$\{x\mid x=k\pi\pm\dfrac{2\pi}{3},\ k\in\mathbb Z\}$　　（D）$\{x\mid x=2k\pi\pm\dfrac\pi3,\ k\in\mathbb Z\}$

\end{question}
\begin{question}
\gkti{11}有一条半径是 $2$ 的弧，其度数是 $60^\circ$，它绕经过弧的中点的直径旋转得到一个球冠，那么这个球冠的面积是\mbox{（　　）}。

（A）$4(2-\sqrt3)\pi$　　（B）$2(2-\sqrt3\pi)$　　（C）$4\sqrt3\pi$　　（D）$2\sqrt3\pi$

\end{question}
\begin{question}
\gkti{12}某小组共有 $10$ 名学生，其中女生 $3$ 名。现选举 $2$ 名代表，至少有 $1$ 名女生当选的不同的选法共有\mbox{（　　）}。

（A）$27$ 种　　（B）$48$ 种　　（C）$21$ 种　　（D）$24$ 种

\end{question}
\begin{question}
\gkti{13}设全集 $U=\mathbb R$，集合 $M=\{x\mid\sqrt{x^2}>2\}$，$N=\{x\mid\log_x7>\log_3 7\}$，那么 $M\cap\overline N=$\mbox{（　　）}。

（A）$\{x\mid x<-2\}$　　（B）$\{x\mid x<-2\text{ \text{或} }x\ge3\}$

（C）$\{x\mid x\ge3\}$　　（D）$\{x\mid -2\le x<3\}$

\end{question}
\begin{question}
\gkti{14}设 $\{a_n\}$ 是由正数组成的等比数列，公比 $q=2$，且 $a_1a_2a_3\cdots a_{30}=2^{30}$，那么 $a_3a_6a_9\cdots a_{30}$ 等于\mbox{（　　）}。

（A）$2^{10}$　　（B）$2^{20}$　　（C）$2^{16}$　　（D）$2^{15}$

\end{question}
\begin{question}
\gkti{15}设 $\triangle ABC$ 不是直角三角形，$A$ 和 $B$ 是它的两个内角，那么\mbox{（　　）}。

（A）“$A<B$”是“$\tan A<\tan B$”的充分条件，但不是必要条件

（B）“$A<B$”是“$\tan A<\tan B$”的必要条件，但不是充分条件

（C）“$A<B$”是“$\tan A<\tan B$”的充分必要条件

（D）“$A<B$”是“$\tan A<\tan B$”的充分条件，也不是必要条件

\end{question}
\begin{question}
\gkti{16}对于定义域是 $\mathbb R$ 的任何奇函数 $f(x)$，都有\mbox{（　　）}。

（A）$f(x)-f(-x)>0\ (x\in\mathbb R)$　　（B）$f(x)-f(-x)\le0\ (x\in\mathbb R)$

（C）$f(x)f(-x)\le0\ (x\in\mathbb R)$　　（D）$f(x)f(-x)>0\ (x\in\mathbb R)$

\end{question}
\begin{question}
\gkti{17}如果双曲线的两条渐近线的方程是 $y=\pm\dfrac32x$，焦点坐标是 $(-\sqrt{26},0)$ 和 $(\sqrt{26},0)$，那么它的两条准线之间的距离是\mbox{（　　）}。

（A）$\dfrac8{13}\sqrt{26}$　　（B）$\dfrac4{13}\sqrt{26}$　　（C）$\dfrac{18}{13}\sqrt{26}$　　（D）$\dfrac9{13}\sqrt{26}$

\end{question}
\begin{question}
\gkti{18}$\tan\dfrac\pi8=\underline{\qquad\qquad}$。

\end{question}
\begin{question}
\gkti{19}设直线的参数方程是 $\begin{cases}x=2+\dfrac12t,\\y=3+\dfrac{\sqrt3}{2}t,\end{cases}$ 那么它的斜截式方程是\underline{\qquad\qquad}。

\end{question}
\begin{question}
\gkti{20}如果三角形的顶点分别是 $O(0,0)$、$A(0,15)$、$B(-8,0)$，那么它的内切圆方程是\underline{\qquad\qquad}。

\end{question}
\begin{question}
\gkti{21}$\displaystyle\lim_{n\to\infty}\left[\frac1{1\cdot4}+\frac1{4\cdot7}+\frac1{7\cdot10}+\cdots+\frac1{(3n-2)(3n+1)}\right]=$\underline{\qquad\qquad}。

\end{question}
\begin{question}
\gkti{22}$91^{92}$ 除以 $100$ 的余数是\underline{\qquad\qquad}。

\end{question}
\begin{question}
\gkti{23}已知三棱锥 $A-BCD$ 的体积是 $V$，棱 $BC$ 的长是 $a$，面 $ABC$ 和面 $DBC$ 的面积分别是 $S_1$ 和 $S_2$。设面 $ABC$ 和面 $DBC$ 所成的二面角是 $\alpha$，那么 $\sin\alpha=$\underline{\qquad\qquad}。

\end{question}
\begin{question}
\gkti{24}已知关于 $x$ 的方程 $2a^{2x-2}-7a^{x-1}+3=0$ 有一个根是 $2$，求 $a$ 的值和方程其余的根。

\end{question}
\begin{question}
\gkti{25}已知平面 $\alpha$ 和不在这个平面内的直线 $a$ 都垂直于平面 $\beta$。求证：$a\parallel\alpha$。

\end{question}
\begin{question}
\gkti{26}证明不等式：$1+\dfrac1{\sqrt2}+\dfrac1{\sqrt3}+\cdots+\dfrac1{\sqrt n}<2\sqrt n\quad(n\in\mathbb N^*)$。

\end{question}
\begin{question}
\gkti{27}设抛物线经过两点 $(-1,6)$ 和 $(-1,-2)$ 对称轴与 $x$ 轴平行，开口向右，直线 $y=2x+7$ 被抛物线截得的线段的长是 $4\sqrt{10}$，求抛物线的方程。

\end{question}
\begin{question}
\gkti{28}求同时满足下列两个条件的所有复数 $z$：

① $z+\dfrac{10}{z}$ 是实数，且 $1<z+\dfrac{10}{z}\le6$；

② $z$ 的实部和虚部都是整数。

\end{question}